\section{Monogenic orders with nontrivial generic residue degree}
\label{sec:monogenic-residue-degree}

The curvature construction also applies to a distinguished defining
polynomial whose generic quotient is a field.  The resulting modulus
is the residue degree of that field: this is exactly the ambiguity in
the valuation of a determinant under a change of generic basis.

\begin{theorem}\label{thm:monogenic-residue-degree}
Let $O=\Z_p$ or $\Z_{(p)}$, and let $F\in O[t]$ be monic of degree
$N\ge2$, with
$F\bmod p=t^N$.  Suppose that
\[
 E=\Q_p[t]/(F)
\]
is a field of residue degree $f>1$ over $\Q_p$.  Put
$A=O[t]/(F)$.  There is a homomorphism
\[
 \Theta:K_0(\Lat_A(S_p))\longrightarrow
             \bigoplus_{i\ge1}\Z/f
\]
which vanishes on the Young span and whose restriction to the kernel
of specialization to $\operatorname{Frac}(O)\otimes_O A$ is
surjective.  In particular $\Phi_A$
fails to be surjective in degree $p$.

The proof supplies explicit generic-zero preimages of the coordinate
vectors.  It asserts neither a torsion order for those preimages nor
a splitting of this quotient.
\end{theorem}

\begin{corollary}\label{cor:scaled-unramified-orders}
Let $E/\Q_p$ be unramified of degree $N\ge2$, choose a unit
$\theta$ with $O_E=\Z_p[\theta]$, and let $a\ge1$.  Every proper
local order $A=\Z_p[p^a\theta]$ satisfies the theorem with modulus
$f=N$.  This includes every degree $N\ge2$ and every prime.
\end{corollary}

\subsection{Curvature for a distinguished defining polynomial}

We first work over $O=\Z_p$; descent of the explicit witnesses to
$\Z_{(p)}$ is proved at the end.  The ring $A$ is finite free and
complete over $O$, and is local with
maximal ideal $(p,t)$.  Formal division by the distinguished polynomial
identifies $A$ with $R/(F)$ for $R=O[[t]]$.  Since its generic quotient
is a field, it is a domain and $F(0)\ne0$.

Let $M$ be an $A$-free $C_p$-lattice, with generator matrix $T$.
Choose a lift $\widetilde T$ over $R$ and put
\[
 \widetilde U=\frac{\widetilde T^p-I}{F},\qquad
 U=\widetilde U\bmod F.
\]
The numerator is divisible by $F$, and $U$ commutes with $T$.
The stable-curvature identity of
Lemma~\ref{lem:truncation-stable-curvature} remains valid with $F$
in place of $t^n$.  Indeed, if an equivariant map $H:M\to M'$ is
lifted and
$\widetilde T'\widetilde H-\widetilde H\widetilde T=F\widetilde B$,
the same telescoping calculation gives
\begin{equation}\label{eq:monogenic-stable-curvature}
 U'H-HU=\sum_{j=0}^{p-1}(T')^{p-1-j}BT^j.
\end{equation}
This is the group trace appearing in that lemma, with the same
factorization through an induced lattice.  Consequently the stable
endomorphism $U$ is natural and independent of lifts and bases.

Use the integral primary decomposition of the finite free $O$-module
$M$ for $U$, indexed by monic irreducibles $q\in\F_p[u]$.
The projectors commute with $A$ and $T$, so their images $M_q$ are
$A$-free direct summands.  The complete-local primary construction and
its exactness from Section~\ref{sec:all-truncations} apply without
change: a common annihilating polynomial gives compatible projectors
throughout a short exact sequence.

We verify the altered acyclicity step explicitly.  For $q\ne u$,
$U$ is a unit on $M_q$.  Choose an $A$-basis adapted to this summand
and lift it to an invertible matrix over $R$.  The off-diagonal blocks
of the lifted action are divisible by $F$.  Every term leaving a
block and returning in its $p$th power contains two such blocks.
Thus the diagonal action lift $\widetilde T_q$ satisfies
\[
 \frac{\widetilde T_q^p-I}{F}\bmod F=U|_{M_q}.
\]
Its lifted curvature $\widetilde U_q$ is a unit over $R$, and
\begin{equation}\label{eq:monogenic-factorization}
 (\widetilde T_q-I)N(\widetilde T_q)=F\widetilde U_q.
\end{equation}
If $(\widetilde T_q-I)x=Fy$, multiplication by the norm, cancellation
of the nonzerodivisor $F$, and inversion of $\widetilde U_q$ give
$x=N(\widetilde T_q)\widetilde U_q^{-1}y$ modulo $F$.
Starting with $N(\widetilde T_q)x=Fy$ gives the other Tate parity.
Hence
\[
 \widehat H^\bullet(C_p,M_q)=0\qquad(q\ne u).
\]
This proves the general-$F$ replacement for
Lemma~\ref{lem:truncation-primary-acyclic}.

\subsection{The determinant ambiguity and the additive coordinates}

For a cohomologically trivial $A$-free cyclic lattice $M$, its invariant
module $M^{C_p}$ is finite free over $O$, and its generic fibre is
an $E$-vector space, say $E^r$.  Choose a generic $E$-basis.  Using
$1,t,\ldots,t^{N-1}$ in each copy of $E$, compare an $O$-basis of
$M^{C_p}$ with the resulting $\Q_p$-basis and define
\[
 \iota_f(M)=v_p(\det M^{C_p})\pmod f.
\]
Integral basis changes have unit determinant.  If $B\in\mathrm{GL}_r(E)$
changes the generic basis, its determinant over $\Q_p$ is
$\operatorname{Norm}_{E/\Q_p}(\det_E B)$, and
\begin{equation}\label{eq:monogenic-norm-valuation}
 v_p\bigl(\operatorname{Norm}_{E/\Q_p}(x)\bigr)=f\,v_E(x).
\end{equation}
Here $v_E$ is normalized by an $E$-uniformizer.  To see the formula,
units preserve the valuation ring, while multiplication by a
uniformizer has cokernel the residue field, of $O$-length $f$.
Determinants and multiplicativity give the formula for every nonzero
$x$.  Thus $\iota_f$ is independent of the generic basis modulo $f$.

This index is additive on admissible sequences of cohomologically
trivial lattices.  Taking invariants is exact because the first term
has vanishing $H^1$; compatible integral and generic bases make the
determinant block triangular.  It is zero on induced projectives:
their invariants are finite projective, hence free, over $A$.
The stabilization argument of
Lemma~\ref{lem:truncation-stable-index} therefore proves invariance
under relative stable isomorphism.  That argument uses only the
locality of $A$, the actual stabilization isomorphism, and the
vanishing of the index on projectives, all of which hold here.

For $q\ne u$ set
\begin{equation}\label{eq:monogenic-primary-index}
 \eta_q(M)=\iota_f(M_q)\in\Z/f.
\end{equation}
We recall the descent to the Grothendieck group to specify exactly
what is reused.  Two choices of curvature agree stably by
\eqref{eq:monogenic-stable-curvature}.  Over the complete local ring
$A$, a common product of their characteristic polynomials gives one
Hensel--Bezout primary projector for both endomorphisms.  The two
projectors agree stably, so their images are stably isomorphic.
The preceding stable invariance proves that \eqref{eq:monogenic-primary-index}
does not depend on choices.  For an admissible sequence, choose an
$A$-splitting and upper triangular action lifts.  Its curvature is
upper triangular, and exact primary extraction gives an admissible
sequence of the $q$-components.  Their Tate acyclicity and index
additivity prove that $\eta_q$ is a homomorphism.

Constants have constant action lifts with zero curvature, so every
$\eta_q$ vanishes on constant $O$-lattices.  Each class has finite
primary support.  These observations include the stable-primary,
exactness, and finite-support assertions needed below.

\subsection{A parameter and its companion lattice}

Let $\tau$ denote the image of $t$ in $E$.  It is integral, and its
reduction is zero because $F\bmod p=t^N$.  Thus $v_E(\tau)\ge1$.
Since its norm is $F(0)$ up to sign,
\[
 v_p(F(0))=f\,v_E(\tau)\ge f\ge2.
\]
The valuation needed below is therefore automatic except for the
additional parameter calculation when $p=2$ and $v_p(F(0))=2$.
In that case $f=2$.

Write $F(0)=p^v b_0$, with $b_0$ a unit.  First suppose that $p$ is
odd, or that $p=2$ and $v\ge3$.  Let $O_s/O$ be the
unramified extension of degree $s$, and prescribe a nonzero residue
$\eta\in\F_{p^s}$.  Choose a unit $y\in O_s$ reducing to
$\overline b_0\eta$, and put
\begin{equation}\label{eq:monogenic-parameter}
 \alpha=1+p^{v-1}y,\qquad c=\frac{\alpha^p-1}{F(0)}.
\end{equation}
The binomial expansion gives
\[
 \alpha^p-1=p^v y\,U_0,\qquad U_0\in1+pO_s.
\]
For odd $p$ this uses $v\ge2$; for $p=2$ the factor is exactly
$U_0=1+2^{v-2}y$, and $v\ge3$ suffices.  Hence $c$ is a unit
with residue $\eta$, and
\[
 \alpha^p=1+cF(0),\qquad v_p(\alpha-1)=v-1.
\]

For the remaining case $p=2$, $F(0)=4b_0$, choose an odd prime
$s\ge3$ and a nonzero trace-zero element $\eta\in\F_{2^s}$.
Its degree is $s$, because the only proper subfield is $\F_2$,
whose only trace-zero element is zero for odd $s$.
The map $y\mapsto y^2+y$ has kernel $\F_2$ and image the
trace-zero hyperplane.  Choose $\overline y$ with
$\overline y^2+\overline y=\eta$ and lift it to $O_s$.
Both $y$ and $1+y$ are units.  Set
\[
 \alpha=1+2y,\qquad c=\frac{y(1+y)}{b_0}.
\]
Then $\alpha^2=1+cF(0)$, $\overline c=\eta$, and both
$\alpha-1=2y$ and $q_\alpha(1)=\alpha+1=2(1+y)$ have
valuation one.  Here $q_\alpha$ is the polynomial defined next.
Thus the same companion and ideal calculations apply.

Put
\[
 G(t)=\frac{F(t)-F(0)}t,\qquad
 q_\alpha(z)=\frac{z^p-\alpha^p}{z-\alpha}.
\]
Let $C_\alpha$ be the companion of $q_\alpha$, let $a$ be its last
coordinate row, and let $b$ be its first coordinate column.  Over
$A_s=A\otimes_O O_s$, take the rank-$p$ lattice with generator
\begin{equation}\label{eq:monogenic-witness}
 T=\begin{pmatrix}
       \alpha&ct\,a\\
       G(t)b&C_\alpha
    \end{pmatrix}.
\end{equation}
The companion cofactor identity gives, over $O_s[[t]]$,
\[
 \det(zI-\widetilde T)
   =(z-\alpha)q_\alpha(z)-ctG(t)
   =z^p-1-cF(t).
\]
Cayley--Hamilton therefore gives the exact identity
\begin{equation}\label{eq:monogenic-exact-curvature}
 \widetilde T^p-I=cF(t)I_p.
\end{equation}
Thus \eqref{eq:monogenic-witness} is a cyclic group action, and its
curvature is scalar multiplication by $c$.  Restriction of coefficients
gives an $A$-free lattice of rank $ps$.

\subsection{The ideal calculation}

We compute its coinvariants; the norm identifies coinvariants with
invariants because the selected lattice is cohomologically trivial.
This is an $A$-linear isomorphism, so the covolume modulo $f$ agrees
with the invariant index above.

Let $a_0=\alpha-1$ and
\[
 \lambda=\frac{q_\alpha(1)}c=\frac{F(0)}{a_0}.
\]
Its valuation is one, so $\lambda$ is $p$ times a unit.  The companion
columns identify all cyclotomic coinvariant coordinates with one
generator $w$.  With $v_0$ the remaining coordinate, the presentation
is
\begin{equation}\label{eq:monogenic-coinvariant-presentation}
 tv_0=\lambda w,\qquad a_0v_0+G(t)w=0.
\end{equation}
The map $v_0\mapsto\lambda$, $w\mapsto t$ identifies this module
with $(p,t)A_s$.  To check completeness of the relations, suppose in
$O_s[[t]]$ that $\lambda a+tb=Fc_0$.  Since
$F=a_0\lambda+tG$, this gives
\[
 \lambda(a-a_0c_0)=t(Gc_0-b).
\]
Setting $t=0$ shows $a-a_0c_0=tq_0$, and then
$b=Gc_0-\lambda q_0$.  Every relation is therefore generated by
$(t,-\lambda)$ and $(a_0,G)$, as in
\eqref{eq:monogenic-coinvariant-presentation}.

Relative to an integral basis of $O_s/O$, the module $A_s$ is free
of rank $s$ over $A$, and
\[
 A_s/(p,t)A_s=\F_{p^s}
\]
has $O$-length $s$.  If $\eta$ has degree $s$ over $\F_p$ and minimal
polynomial $q$, the whole witness lies in its $q$-primary component.
Consequently
\begin{equation}\label{eq:monogenic-witness-index}
 \eta_q(L)=s\pmod f,
\end{equation}
and its other primary coordinates vanish.  This computation does not
require $E\otimes_{\Q_p}\operatorname{Frac}(O_s)$ to be a field.

\subsection{The countable quotient and the generic kernel}

Outside the valuation-two case at $p=2$, choose distinct positive
degrees $s_i$ coprime to $f(p-1)$, and monic
irreducibles $q_i\in\F_p[u]$ of those degrees with nonzero constant
terms.  Such irreducibles exist in every degree.  Choose a root
$\eta_i$ of each.  In the valuation-two case, use trace-zero
parameters of distinct odd prime degrees from the preceding
construction; these degrees are prime to $f=2$.
Form the corresponding coefficient-restricted lattices $L_i$ and put
\[
 x_i=[\operatorname{Ind}_{C_p}^{S_p}L_i]-s_i[A[S_p]].
\]
Under $g\mapsto g^j$, the lift $\widetilde T^j$ has residual
curvature $j\eta_i$.  If this scalar twist preserves $q_i$, comparison
of constant terms gives $j^{s_i}=1$, so $j=1$ by the degree condition.
Different degrees have disjoint scalar orbits.

Mackey restriction contributes each of the $p-1$ normalizer twists
once.  Every other intersection is trivial, and gives a relatively
projective cyclic lattice.  Therefore
\[
 (\eta_{q_j}\circ\operatorname{Res}_{C_p}^{S_p})(x_i)
       =s_i\delta_{ij}\pmod f.
\]
Normalize the $i$th coordinate by $s_i^{-1}$ modulo $f$.  Finite
primary support gives the homomorphism $\Theta$ to the direct sum,
and these equations show surjectivity.  Every Young lattice is
constant over $O$, so $\Theta$ kills the Young span.

To verify the asserted generic vanishing, work over
$E_s=E\otimes_{\Q_p}\operatorname{Frac}(O_s)$.  The elements $t$
and $c$ are units.  Starting with the first cyclotomic coordinate in
\eqref{eq:monogenic-witness}, its successive images first span all
the cyclotomic coordinates, and the next image has remaining
coordinate $ct$, a unit.  They form a cyclic basis.  Since the
characteristic polynomial is $z^p-1$, the generic lattice is
$E_s[C_p]$.  Restriction of coefficients and induction show that
$x_i$ is generically zero.  Thus $\Theta$ is surjective already on
the generic kernel in the complete case.

\subsection{Descent to the localized integers}

Now let $O=\Z_{(p)}$.  The hypothesis that $\Q_p[t]/F$ is a field
implies that $F$ is irreducible over $\Q$.  The ring $A$ is finite
free and local over $O$, with maximal ideal $(p,t)$, and its
completion is $\widehat A=\Z_p[t]/F$.  Compose the complete
coordinate homomorphisms with coefficient extension to $\widehat A$.
Admissible sequences remain exact because they split over $A$.
This gives finite-support homomorphisms on the original lattice
group, vanishing on its Young span.

Every required witness is already defined over $O$.  Lift a monic
irreducible polynomial for $\F_{p^s}/\F_p$ to $O[z]$ and take its
quotient $O_s$.  This is a finite etale local $O$-algebra: its
residue algebra is the separable field $\F_{p^s}$, and its maximal
ideal is $pO_s$.  It is therefore a DVR, free of rank $s$ over $O$,
and its completion is the unramified extension used above.
Lift each chosen residue $y$ to this ring.  Formulas
\eqref{eq:monogenic-parameter}, or the valuation-two formula
$c=y(1+y)/b_0$, are finite identities in $O_s$; the valuation
calculations prove that their quotients are units there.
No Hensel root in $\Z_{(p)}$ is needed.

The resulting lattices complete to witnesses with the computed
coordinate vectors, so the pulled-back map is surjective.
Their generic cancellation holds over $\Q$ itself.  Indeed
$A\otimes_O\Q$ is a number field, and after adjoining
$\operatorname{Frac}(O_s)$ it is a product of characteristic-zero
fields.  In every factor $t$ and $c$ are units, and the explicit
cyclic basis above identifies the witness with the regular
representation.  Restriction and induction give generic-zero
preimages over the original ring.  This proves the theorem over
$\Z_{(p)}$ without assuming that arbitrary completed lattices descend.

For Corollary~\ref{cor:scaled-unramified-orders}, let $H$ be the
minimal polynomial of $\theta$.  It has degree $N$ and unit constant
term.  The minimal polynomial of $p^a\theta$ is
\[
 F(t)=p^{aN}H(t/p^a).
\]
It is monic integral, reduces to $t^N$, and has constant-term
valuation $aN$.  The field has residue degree $N$, so all hypotheses
hold.  The order is proper because its residue
field is $\F_p$, whereas that of the maximal order is $\F_{p^N}$.

\begin{example}\label{ex:unramified-cubic-order}
The polynomial $H(z)=z^3-z-1$ has no root over $\F_3$, so it defines
an unramified cubic extension of $\Q_3$.  If $\theta$ is a root, the
order
\[
 A=\Z_3[3\theta]
   \simeq\Z_3[t]/(t^3-9t-27)
\]
has a degree-three generic kernel mapping onto
$\bigoplus_{i\ge1}\Z/3$.  One explicit cyclic witness is obtained
from $\alpha=10$ and $c=-37$:
\[
 T=\begin{pmatrix}
       10&0&-37t\\
       t^2-9&0&-100\\
       0&1&-10
    \end{pmatrix},\qquad T^3=I.
\]
Its induced class minus the regular class has coordinate one modulo
three for the primary label $u-2$.  The theorem supplies independent
coordinates through the unramified parameter extensions; it does not
assign an exact order to this displayed class.
\end{example}
